\documentclass[10pt,a4paper]{amsart}
\usepackage[margin=19mm]{geometry}
\usepackage{amsmath,amssymb,mathtools}
\usepackage[scr=rsfs,cal=euler]{mathalfa}
\usepackage{xfrac}
\usepackage[unicode,hidelinks,bookmarks=false]{hyperref}
\newcommand{\ie}{i.e.\ }
\newcommand{\cf}{cf.\ }
\newcommand{\resp}{respectively\ }
\newcommand{\Subsection}{\subsection}
\providecommand{\nobreakdash}{\nobreak\hbox{-}}
\setlength{\emergencystretch}{2em}
\multlinegap0pt
\usepackage{amssymb}
\usepackage{mathtools}
\usepackage{xcolor}
\newtheorem{theorem}{Theorem}
\newtheorem{corollary}{Corollary}
\usepackage{cleveref}
\crefname{theorem}{Theorem}{Theorems}
\crefname{corollary}{Corollary}{Corollarys}
\newcommand{\Dfun}[2]{D(#1,#2)}
\newcommand{\R}{\mathbb R}
\newcommand{\abs}[1]{\left\lvert #1\right\rvert}
\newcommand{\norm}[1]{\left\lVert #1\right\rVert_2}
\newcommand{\rank}{\operatorname{rank}}
\newcommand{\spann}{\operatorname{span}}
\title{The Sharp Dimension Bound in the Johnson--Lindenstrauss Lemma}
\author{Vishesh Jain}
\date{Source: arXiv:2608.13782v2 (CC BY 4.0)}
\begin{document}
\maketitle

\noindent\emph{Source and licence.} The Sharp Dimension Bound in the Johnson--Lindenstrauss Lemma. Version arXiv:2608.13782v2 (CC BY 4.0), distributed by arXiv under the Creative Commons Attribution 4.0 International licence, http://creativecommons.org/licenses/by/4.0/, as recorded in the arXiv licence metadata for this exact version and shown on the abstract page https://arxiv.org/abs/2608.13782v2. Full licence notice: \texttt{licenses/arxiv-2608.13782-CC-BY-4.0.txt}.\par\smallskip

\noindent\textit{Editorial scope.} Counted unit: the article's theorem thm:psd (source0.tex lines 290--305). The article's complete original proof of it is delivered with it (source0.tex lines 672--793, one \texttt{proof} environment of 122 lines, which the article places away from the statement). No mathematics is written, repaired or re-derived: the statement and the proof are the article's own lines, in the article's own notation, and no retained line is substituted or re-flowed.  Retained uncounted as the source-local context the proof needs: lines 194--207; lines 623--641.  Nothing was omitted from any retained span, so the package carries no omission record.  No retained line contains a citation, so the wrapper carries no bibliography and no bibliography entry.  No retained line contains a figure environment, an image-inclusion command or drawing code, so no image is delivered and the package declares no asset dependency.  Editorial formatting only, in addition: an AMS \texttt{amsart} preamble that restates the article's own declarations and macros used by the retained lines, namely the environments \texttt{theorem}, \texttt{corollary} on separate counters, together with the standard packages those lines need.  The retained environments print in this document as Theorem 1, Corollary 1 in order of appearance, whereas the article prints the same items with the numbering of its own full document.  The complete native source of the article as distributed in the arXiv e-print for this exact version is bundled unchanged as \texttt{sources/207613/source0.tex}.
\begin{theorem}
\label{thm:jl-main}
There is a universal constant $C>0$ such that the following holds. Let
$n,d\geq2$, let $0<\varepsilon<1/2$, and let $X\subseteq\R^d$ be a set
of $n$ points. Then there is an integer
\[
 r\leq
 C\min\left\{
 d,n-1,\frac{\log(2+\varepsilon^2n)}{\varepsilon^2}
 \right\}
\]
and a linear map $L:(\R^d, \|\cdot \|_2)\to (\R^r, \|\cdot \|_2)$ whose restriction to $X$ has distortion
at most $1+\varepsilon$. 
\end{theorem}

\begin{corollary}\label{cor:one-step}
There is a universal constant $C_0>0$ such that the following holds.
Let $n\geq2$ and $r\geq1$. For every finite set
$X\subseteq\R^{2r}$ with $2\leq\abs{X}\leq n$, there is a linear map
$L:\R^{2r}\to\R^r$ such that
\[
 (1-\delta_r)\norm{x-y}^2
 \leq
 \norm{Lx-Ly}^2
 \leq
 (1+\delta_r)\norm{x-y}^2
\]
for all $x,y\in X$, where
\[
 \delta_r
 =
 C_0\sqrt{\frac{\log(2+n/r)}{r}}.
\]
\end{corollary}

\begin{theorem}
\label{thm:psd}
Let $r,m\geq1$ and let $u_1,\ldots,u_m\in\R^{2r}$ be unit vectors. There is
a matrix $M\succeq0$ such that
\[
 \rank M\leq r
\]
and
\[
 \max_{\ell\leq m}
 \abs{u_\ell^{\mathsf T}Mu_\ell-1}
 \leq
 C\sqrt{\frac{\log(2+m/r^2)}{r}},
\]
where $C>0$ is a universal constant.
\end{theorem}

\begin{proof}[Proof of \Cref{thm:jl-main} given \Cref{thm:psd}]
Translate the configuration by $-x_1$ and write
\[
 \widetilde X=\{x-x_1:x\in X\},
 \qquad
 V=\spann(\widetilde X).
\]
Translation does not change pairwise distances, and
\[
 \dim V\leq\min\{d,n-1\}.
\]
It suffices to construct the desired linear map on $V$, since the map
may then be extended by zero on $V^\perp$.

Let $C_0$ be the constant in \Cref{cor:one-step}, and choose a
sufficiently large universal constant $K$. Set
\[
 r_0=
 \left\lceil
 K\Dfun{n}{\varepsilon}
 \right\rceil.
\]
If $\dim V\leq r_0$, an isometric identification of $V$ with
$\R^{\dim V}$ gives an exact realization. Its target dimension satisfies
\[
 \dim V\leq\min\{d,n-1,r_0\}
 \leq
 C\min\{d,n-1,\Dfun{n}{\varepsilon}\}.
\]
We may therefore assume that $\dim V>r_0$.

Choose $s\geq1$ minimal such that
\[
 \dim V\leq2^sr_0.
\]
For $0\leq j\leq s-1$, set
\[
 \beta_j
 =
 C_0\sqrt{
 \frac{\log(2+n/(2^jr_0))}{2^jr_0}
 },
 \qquad
 S=\sum_{j=0}^{s-1}\beta_j.
\]
Since $j\mapsto\log(2+n/(2^jr_0))$ is decreasing,
\[
 S
 \leq
 C_0\sum_{j=0}^{\infty}2^{-j/2}
 \sqrt{\frac{\log(2+n/r_0)}{r_0}}
 \leq \frac{C\varepsilon}{\sqrt K}.
\]
Indeed, for $K$ larger than a universal constant,
\[
 r_0
 \geq K\varepsilon^{-2}\log(2+\varepsilon^2n)
 \geq \varepsilon^{-2},
\]
and hence
\[
 \log(2+n/r_0)\leq\log(2+\varepsilon^2n).
\]
Choose $K$ sufficiently large that $S\leq\varepsilon/10$. In
particular, every $\beta_j$ lies in $[0,1)$.

Fix an isometric linear embedding
\[
 \iota:V\longrightarrow\R^{2^sr_0}
\]
and set $X_s=\iota(\widetilde X)$. For
$j=s-1,s-2,\ldots,0$, apply \Cref{cor:one-step} to $X_{j+1}$, with
target dimension $2^jr_0$, obtaining a linear map
\[
 L_j:\R^{2^{j+1}r_0}\longrightarrow\R^{2^jr_0},
 \qquad
 X_j=L_j(X_{j+1}).
\]
The squared-distance error at this step is at most $\beta_j$. Note that since
$\beta_j<1$, the map $L_j$ is injective on $X_{j+1}$.

Let
\[
 A=L_0\circ L_1\circ\cdots\circ L_{s-1}\circ\iota.
\]
For every distinct $x,y\in X$, the ratio between the final and initial
squared distances lies between
\[
 \prod_{j=0}^{s-1}(1-\beta_j)
 \quad\text{and}\quad
 \prod_{j=0}^{s-1}(1+\beta_j).
\]
Consequently,
\[
 \prod_{j=0}^{s-1}(1-\beta_j)\geq1-S,
 \qquad
 \prod_{j=0}^{s-1}(1+\beta_j)\leq e^S.
\]
The distortion of $A$ on $\widetilde X$ is therefore at most
\[
 \sqrt{\frac{e^S}{1-S}}
 \leq e^{3S/2}
 \leq1+\varepsilon,
\]
where we used $S\leq\varepsilon/10$ and
$0<\varepsilon<1/2$.

The target dimension of this embedding is $r_0$.  Since $r_0<\dim V$ in the case under
consideration and $r_0\leq C\Dfun{n}{\varepsilon}$,
\[
 r_0\leq
 C\min\left\{
 d,n-1,\Dfun{n}{\varepsilon}
 \right\}
\]
for a universal constant $C$.  Finally, extend $A$ by zero on
$V^\perp$.  For $x,y\in X$, this extension sends $x-y$ to
\[
 A(x-y)=A\bigl((x-x_1)-(y-x_1)\bigr),
\]
so it has the same distortion on the original configuration $X$.
\end{proof}

\end{document}
